\SourceLeaf{322}
\[
\begin{array}{rcll@{\qquad}rcll}
2 & \text{by} & 180^0+A & \text{to the right of }1 & 3 & \text{by} & B & \text{to the right of }1\\[2pt]
 & \text{,,} & A & \text{,,} & & \text{,,} & 180^0+B & \text{,,}\\[2pt]
 & \text{,,} & 180^0+A & \text{,,} & & \text{,,} & 180^0+B & \text{,,}\;.
\end{array}
\]

The result itself, however, always agrees with only one of these four sets of data, and the other three are \textls[120]{physically impossible}.

The metaphysical reason for this phenomenon is that all one uses of the observed direction is the fact that the \textls[120]{straight lines} make certain angles with each other. In doing so one allows these straight lines to extend in indefinitum on \textls[120]{both} sides, whereas light advances in only one definite sense, so that the cases in which the object would lie in the reverse direction have to be excluded.

For completeness, one must therefore solve the problem of telling at once from the data, before carrying out the calculation, whether they are compatible. This can be done very neatly, but I do not wish to forestall your working it out yourself. \mbox{. . . . . . .}



\Needspace{5\baselineskip}
\centerline{\textsc{Gauss} to \textsc{Gerling}. Göttingen, 14 January 1842.}

\indent \mbox{. . . . . . . .} I am happy to communicate the criterion for \textsc{Pothenot}'s problem to you, but I would like you to keep it to yourself. My reason is that I would like to reserve the treatment of the associated theorem for myself on a suitable occasion, less on account of the elegance of the theorem in itself than on account of the elegance afforded here by the application of complex quantities. In particular, I would want an occasion on which I can say more about the use of complex quantities. If the theorem were made public beforehand from another quarter, I would no longer feel inclined afterwards to return to it myself. Now to the matter at hand.

The criterion can be cast in many forms. The neatest way will be to present it so that all three given points enter symmetrically. I therefore denote the three given points by 1, 2, 3;\space%
